% Generated draft: requires mathematical and source review.
\documentclass[11pt]{article}
\usepackage[margin=25mm]{geometry}
\usepackage{fontspec}
\setmainfont{Latin Modern Roman}
\usepackage{amsmath,amssymb,mathtools,unicode-math}
\setmathfont{Latin Modern Math}
\usepackage{enumitem,needspace,xurl,hyperref}
\setlength{\parindent}{0pt}
\setlength{\parskip}{5pt}
\setlength{\emergencystretch}{3em}
\newcommand{\N}{\mathbb{N}}
\newcommand{\Z}{\mathbb{Z}}
\newcommand{\Q}{\mathbb{Q}}
\newcommand{\R}{\mathbb{R}}
\newcommand{\C}{\mathbb{C}}
\newcommand{\F}{\mathbb{F}}
\newcommand{\A}{\mathbb{A}}
\newcommand{\PP}{\mathbb{P}}
\newcommand{\E}{\mathbb{E}}
\newcommand{\eps}{\varepsilon}
\newcommand{\dd}{\,\mathrm{d}}
\begin{document}
\section*{UnsolvedMath ID: 30003411}
\textbf{Source problem code:} catalogue entry
\subsection{Definitions and mathematical statement}
A tree is a partially ordered set \((T,<_{T})\) such that the strict predecessors of each node are well ordered by \(<_{T}\). The height of a node is the order type of its predecessors, and \(T_{\xi}\) denotes the set of nodes of height \(\xi\). For an infinite cardinal \(\lambda\), a \(\lambda\)-Aronszajn tree is a tree of height \(\lambda\) such that \(|T_{\xi}|<\lambda\) for every \(\xi<\lambda\), with no chain of order type \(\lambda\). For an infinite cardinal \(\mu\), a \(\mu^{+}\)-Aronszajn tree is special if there is a function \(f:T\to\mu\) such that \(f(s)\ne f(t)\) whenever \(s<_{T}t\); equivalently, its nodes are the union of at most \(\mu\) antichains.

For a limit ordinal \(\kappa\), a set \(C\subseteq\kappa\) is closed unbounded, or a club, if it is unbounded in \(\kappa\) and contains every nonzero limit ordinal \(\delta<\kappa\) satisfying \(\sup(C\cap\delta)=\delta\).

The primary question is whether there is a Radin forcing \(\mathbb P\), at an appropriate \(\kappa\), which adds a club \(C\subseteq\kappa\) such that in the forcing extension
\[
\forall\mu\in C\;\bigl(\mu\text{ is a cardinal}\ \Longrightarrow\ \text{there is no special }\mu^{+}\text{-Aronszajn tree}\bigr).
\]
Here \(\mu^{+}\) denotes the successor cardinal computed in the extension, under the usual interpretation of a forcing assertion. The quantifier ranges over cardinal members of \(C\), and the assertion must hold simultaneously for all such members. The source supplies neither the hypotheses on \(\kappa\) nor the measure sequence or precise variant of Radin forcing. Thus this is a schematic existence question, not a fully specified forcing construction or relative-consistency assertion. The tree definitions above concern infinite cardinals; the source does not explicitly restrict the cardinal members of \(C\) to be infinite.

The original statement also asks a second question, absent from the cleaned statement: is it consistent that
\[
\forall\alpha<\omega_{1}\;\bigl(\alpha\text{ is a limit ordinal}\ \Longrightarrow\ \text{there is no }\aleph_{\alpha+1}\text{-Aronszajn tree}\bigr)?
\]
The ordinals and cardinals in this assertion are computed in the proposed model. The source does not specify the base theory or consistency-strength assumptions, or whether its convention for limit ordinals includes \(0\). These qualifications remain unresolved. This second assertion excludes all Aronszajn trees at the indicated cardinals; the first excludes only special Aronszajn trees. No equivalence or implication between the two questions is supplied.
\subsection{Short English statement}
The primary question asks whether Radin forcing can add a club whose cardinal members simultaneously have no special Aronszajn trees at their successors. The original record additionally asks for the consistency of having no Aronszajn trees at \(\aleph_{\alpha+1}\) for every countable limit ordinal \(\alpha\). The forcing hypotheses, consistency assumptions, and convention concerning \(\alpha=0\) require clarification.
\subsection{Sources}
\begin{itemize}
\item Local UnsolvedMath registry: \texttt{data/unsolvedmath/problems.json}, permanent ID \(30003411\), source problem code \texttt{catalogue entry}, title ``Radin Forcing and Simultaneous Tree Properties.'' The supplied record includes both the cleaned primary question and an original statement containing a second consistency question. No individual proposer is identified; the proposed year is recorded as 2017.
\item Supplied report citation: \emph{Set Theory (2018)}, Oberwolfach Reports, pp.\ 527--589, DOI \texttt{10.4171/owr/2017/11}: \url{https://doi.org/10.4171/owr/2017/11}. The record attributes the problem to workshop open-problem discussions documented in Oberwolfach Reports.
\item Supplied literature source, ``OWR Set Theory report'': \url{https://ems.press/content/serial-article-files/46674}.
\item Supplied literature source, ``Diagonal supercompact Radin forcing'': \url{https://doi.org/10.1016/j.apal.2020.102828}. No authors or further bibliographic details are supplied.
\item Literature checked at: \texttt{2026-08-21}. This exact date identifies supplied dated curation, not an independently verified assessment. The supplied assessment labels the question open, says the cited work records it as open, and states that later global forcing results remain partial and do not establish the exact simultaneous assertion. The background explicitly requests independent verification of time-sensitive status and cited claims.
\item The registry labels the statement \texttt{corrected\_verified} and reports that the official Oberwolfach report was checked directly. These are supplied provenance claims; this draft does not independently verify the report, the formulation, or the open-status label.
\end{itemize}
\end{document}
